% GENERATED FILE. Edit canonical lessons, then run npm run generate:books.
% canonical-source-hash: fnv1a64:6507840e0c7dd981
% canonical-lessons: ML-C517-vikaram, ML-C517-ahladam, ML-C517-alaruka, ML-C517-kopam, ML-C517-visamikkuka

\chapter{Feelings Run High}
\label{ch:ml-feelings-run-high}
\hlchaptermodality{517}

\begin{chapteropening}
\textbf{By the end of this chapter:} \emph{I can say an emotion, delight, to bellow, wrath and to be upset.}

Complete the last lesson of chapter 517: I can say an emotion, delight, to bellow, wrath and to be upset.
\end{chapteropening}

\section[\texorpdfstring{\ml{വികാരം} (vikāraṁ)}{വികാരം (vikāraṁ)}]{\ml{വികാരം} (vikāraṁ) — an emotion; a feeling}

\label{lesson:ML-C517-vikaram}



\begin{quote}
Before the new one: say the Malayalam for a principal, then the Malayalam for tuition.
\end{quote}

\subsection*{You'll want to know: \ml{വികാരം}}
\textbf{\ml{വികാരം}} — \emph{vikāraṁ} — \textquotedblleft{}an emotion; a feeling\textquotedblright{}.

\begin{quote}
\ml{ആ} \ml{പാട്ട്} \ml{എന്നിൽ} \ml{ഒരു} \ml{വലിയ} \ml{വികാരം} \ml{ഉണ്ടാക്കി}.
\end{quote}

\begin{quote}
\emph{ā pāṭṭŭ ennil oru valiya vikāraṁ uṇṭākki.}
\end{quote}

\begin{quote}
That song stirred a strong feeling in me.
\end{quote}

Plural: \ml{വികാരങ്ങൾ}, \emph{vikāraṅṅaḷ}.

Plural: \ml{വികാരങ്ങൾ}. A word for any feeling, mild or strong; \ml{വികാരം} \ml{കൊള്ളുക} is to be carried away by one.

\begin{cousinweb}
\textbf{\ml{വികാരം}} (\emph{vikāraṁ}) is borrowed from Sanskrit \emph{vikāra}, \textquotedblleft{}a change, a disturbance of body or mind, a passion\textquotedblright{} (CDIAL 11624). The sense of an emotion grew from the idea of the mind being moved out of its usual state.
\end{cousinweb}

\begin{grammarlens}[title={What you've built}]
One more word for describing what you lived and felt.
\end{grammarlens}

\subsection*{Your turn}
\begin{itemize}
\raggedright
  \item \emph{Say it:} \emph{vikāraṁ}
  \item \emph{Say it:} \emph{vikāraṁ}, once more
  \item \emph{Say it:} \emph{ā pāṭṭŭ ennil oru valiya vikāraṁ uṇṭākki.} — That song stirred a strong feeling in me.
  \item \emph{Recall:} say the Malayalam for a principal, then the Malayalam for tuition, then say \emph{vikāraṁ} again
\end{itemize}

\begin{tcolorbox}[breakable,colback=teal!4,colframe=teal!35!black,title={Before you move on}]
What does \ml{വികാരം} mean? (\textquotedblleft{}An emotion; a feeling\textquotedblright{}.) Say it once more.
\end{tcolorbox}
\section[\texorpdfstring{\ml{ആഹ്ലാദം} (āhlādaṁ)}{ആഹ്ലാദം (āhlādaṁ)}]{\ml{ആഹ്ലാദം} (āhlādaṁ) — delight; glee}

\label{lesson:ML-C517-ahladam}



\begin{quote}
Before the new one: say the Malayalam for tuition, then the Malayalam for an emotion.
\end{quote}

\subsection*{You'll want to know: \ml{ആഹ്ലാദം}}
\textbf{\ml{ആഹ്ലാദം}} — \emph{āhlādaṁ} — \textquotedblleft{}delight; glee\textquotedblright{}.

\begin{quote}
\ml{പരീക്ഷയിൽ} \ml{ജയിച്ചപ്പോൾ} \ml{കുട്ടികൾ} \ml{ആഹ്ലാദത്തോടെ} \ml{ചാടി}.
\end{quote}

\begin{quote}
\emph{parīkṣayil jayiccappōḷ kuṭṭikaḷ āhlādattōṭe cāṭi.}
\end{quote}

\begin{quote}
When they passed the exam, the children jumped for joy.
\end{quote}

With -\ml{ഓടെ} it makes a manner phrase: \ml{ആഹ്ലാദത്തോടെ}, with delight. A -\ml{ം} noun takes -\ml{ത്ത}- before an ending.

\begin{cousinweb}
\textbf{\ml{ആഹ്ലാദം}} (\emph{āhlādaṁ}) is borrowed from Sanskrit \emph{āhlāda}, \textquotedblleft{}refreshment, joy, delight\textquotedblright{} (Monier-Williams). Like many Sanskrit nouns in -a, it ends in -\ml{ം} in Malayalam.
\end{cousinweb}

\begin{grammarlens}[title={What you've built}]
One more word for describing what you lived and felt.
\end{grammarlens}

\subsection*{Your turn}
\begin{itemize}
\raggedright
  \item \emph{Say it:} \emph{āhlādaṁ}
  \item \emph{Say it:} \emph{āhlādaṁ}, once more
  \item \emph{Say it:} \emph{parīkṣayil jayiccappōḷ kuṭṭikaḷ āhlādattōṭe cāṭi.} — When they passed the exam, the children jumped for joy.
  \item \emph{Recall:} say the Malayalam for tuition, then the Malayalam for an emotion, then say \emph{āhlādaṁ} again
\end{itemize}

\begin{tcolorbox}[breakable,colback=teal!4,colframe=teal!35!black,title={Before you move on}]
What does \ml{ആഹ്ലാദം} mean? (\textquotedblleft{}Delight; glee\textquotedblright{}.) Say it once more.
\end{tcolorbox}
\section[\texorpdfstring{\ml{അലറുക} (alaṟuka)}{അലറുക (alaṟuka)}]{\ml{അലറുക} (alaṟuka) — to bellow; to roar}

\label{lesson:ML-C517-alaruka}



\begin{quote}
Before the new one: say the Malayalam for an emotion, then the Malayalam for delight.
\end{quote}

\subsection*{You'll want to know: \ml{അലറുക}}
\textbf{\ml{അലറുക}} — \emph{alaṟuka} — \textquotedblleft{}to bellow; to roar\textquotedblright{}.

\begin{quote}
\ml{ദേഷ്യം} \ml{വന്നപ്പോൾ} \ml{അയാൾ} \ml{അലറി}.
\end{quote}

\begin{quote}
\emph{dēṣyaṁ vannappōḷ ayāḷ alaṟi.}
\end{quote}

\begin{quote}
When he got angry, he bellowed.
\end{quote}

Past: \ml{അലറി}, \emph{alaṟi}.

Said of an angry person, a tiger or a storm; the past is \ml{അലറി}.

\begin{cousinweb}
\textbf{\ml{അലറുക}} (\emph{alaṟuka}) is native Dravidian: DEDR 245 gives the Malayalam verb as to roar or bellow, as an elephant or a tiger does (DEDR 245).
\end{cousinweb}

\begin{grammarlens}[title={What you've built}]
One more word for describing what you lived and felt.
\end{grammarlens}

\subsection*{Your turn}
\begin{itemize}
\raggedright
  \item \emph{Say it:} \emph{alaṟuka}
  \item \emph{Say it:} \emph{alaṟuka}, once more
  \item \emph{Say it:} \emph{dēṣyaṁ vannappōḷ ayāḷ alaṟi.} — When he got angry, he bellowed.
  \item \emph{Recall:} say the Malayalam for an emotion, then the Malayalam for delight, then say \emph{alaṟuka} again
\end{itemize}

\begin{tcolorbox}[breakable,colback=teal!4,colframe=teal!35!black,title={Before you move on}]
What does \ml{അലറുക} mean? (\textquotedblleft{}To bellow; to roar\textquotedblright{}.) Say it once more.
\end{tcolorbox}
\section[\texorpdfstring{\ml{കോപം} (kōpaṁ)}{കോപം (kōpaṁ)}]{\ml{കോപം} (kōpaṁ) — wrath; anger}

\label{lesson:ML-C517-kopam}



\begin{quote}
Before the new one: say the Malayalam for delight, then the Malayalam for to bellow.
\end{quote}

\subsection*{You'll want to know: \ml{കോപം}}
\textbf{\ml{കോപം}} — \emph{kōpaṁ} — \textquotedblleft{}wrath; anger\textquotedblright{}.

\begin{quote}
\ml{അച്ഛന്റെ} \ml{കോപം} \ml{പെട്ടെന്ന്} \ml{മാറി}.
\end{quote}

\begin{quote}
\emph{acchanṟe kōpaṁ peṭṭennŭ māṟi.}
\end{quote}

\begin{quote}
Father's anger passed quickly.
\end{quote}

More formal than \ml{ദേഷ്യം}, and common in stories and old songs.

\begin{cousinweb}
\textbf{\ml{കോപം}} (\emph{kōpaṁ}) is borrowed from Sanskrit \emph{kopa}, \textquotedblleft{}fury, anger\textquotedblright{}, from the root \emph{kup}, \textquotedblleft{}to be agitated, to be angry\textquotedblright{} (Monier-Williams).
\end{cousinweb}

\begin{grammarlens}[title={What you've built}]
One more word for describing what you lived and felt.
\end{grammarlens}

\subsection*{Your turn}
\begin{itemize}
\raggedright
  \item \emph{Say it:} \emph{kōpaṁ}
  \item \emph{Say it:} \emph{kōpaṁ}, once more
  \item \emph{Say it:} \emph{acchanṟe kōpaṁ peṭṭennŭ māṟi.} — Father's anger passed quickly.
  \item \emph{Recall:} say the Malayalam for delight, then the Malayalam for to bellow, then say \emph{kōpaṁ} again
\end{itemize}

\begin{tcolorbox}[breakable,colback=teal!4,colframe=teal!35!black,title={Before you move on}]
What does \ml{കോപം} mean? (\textquotedblleft{}Wrath; anger\textquotedblright{}.) Say it once more.
\end{tcolorbox}
\section[\texorpdfstring{\ml{വിഷമിക്കുക} (viṣamikkuka)}{വിഷമിക്കുക (viṣamikkuka)}]{\ml{വിഷമിക്കുക} (viṣamikkuka) — to be upset; to feel bad}

\label{lesson:ML-C517-visamikkuka}



\begin{quote}
Before the new one: say the Malayalam for to bellow, then the Malayalam for wrath.
\end{quote}

\subsection*{You'll want to know: \ml{വിഷമിക്കുക}}
\textbf{\ml{വിഷമിക്കുക}} — \emph{viṣamikkuka} — \textquotedblleft{}to be upset; to feel bad\textquotedblright{}.

\begin{quote}
\ml{ജോലി} \ml{പോയതിൽ} \ml{അവൻ} \ml{വളരെ} \ml{വിഷമിച്ചു}.
\end{quote}

\begin{quote}
\emph{jōli pōyatil avan vaḷare viṣamiccu.}
\end{quote}

\begin{quote}
He was very upset at losing his job.
\end{quote}

Past: \ml{വിഷമിച്ചു}, \emph{viṣamiccu}.

The noun is \ml{വിഷമം}, distress. What upsets you takes -\ml{ഇൽ}: \ml{ജോലി} \ml{പോയതിൽ}, at the job going.

\begin{cousinweb}
\textbf{\ml{വിഷമിക്കുക}} (\emph{viṣamikkuka}) is built from Sanskrit \emph{viṣama}, \textquotedblleft{}uneven, difficult, dangerous\textquotedblright{} (CDIAL 11972), with the Malayalam verb ending -\ml{ിക്കുക}. To be upset is to be in an uneven, difficult state.
\end{cousinweb}

\begin{grammarlens}[title={What you've built}]
One more word for describing what you lived and felt.
\end{grammarlens}

\subsection*{Your turn}
\begin{itemize}
\raggedright
  \item \emph{Say it:} \emph{viṣamikkuka}
  \item \emph{Say it:} \emph{viṣamikkuka}, once more
  \item \emph{Say it:} \emph{jōli pōyatil avan vaḷare viṣamiccu.} — He was very upset at losing his job.
  \item \emph{Recall:} say the Malayalam for to bellow, then the Malayalam for wrath, then say \emph{viṣamikkuka} again
\end{itemize}

\begin{tcolorbox}[breakable,colback=teal!4,colframe=teal!35!black,title={Before you move on}]
What does \ml{വിഷമിക്കുക} mean? (\textquotedblleft{}To be upset; to feel bad\textquotedblright{}.) Say it once more.
\end{tcolorbox}
